\documentclass[11pt,a4paper]{article}
\usepackage[margin=26mm]{geometry}
\usepackage[T1]{fontenc}
\usepackage[utf8]{inputenc}
\usepackage{lmodern,microtype}
\usepackage{amsmath,amssymb,amsthm,mathtools,mathrsfs}
\usepackage{booktabs,array,longtable,enumitem,xcolor}
\usepackage{listings,fancyhdr,hyperref}
\definecolor{ink}{RGB}{24,51,76}
\definecolor{accent}{RGB}{0,119,139}
\hypersetup{colorlinks=true,linkcolor=ink,citecolor=accent,urlcolor=accent,
 pdftitle={Toeplitz Counterexamples and Mertens Mean Oscillations},
 pdfauthor={Research prepared for Keio}}
\pagestyle{fancy}\fancyhf{}
\fancyhead[L]{\small\color{ink}Toeplitz counterexamples and Mertens means}
\fancyhead[R]{\small Research note}
\fancyfoot[C]{\thepage}
\setlength{\headheight}{15pt}
\setlength{\emergencystretch}{2em}
\numberwithin{equation}{section}
\newtheorem{theorem}{Theorem}[section]
\newtheorem{proposition}[theorem]{Proposition}
\newtheorem{lemma}[theorem]{Lemma}
\newtheorem{corollary}[theorem]{Corollary}
\theoremstyle{definition}\newtheorem{definition}[theorem]{Definition}
\theoremstyle{remark}\newtheorem{remark}[theorem]{Remark}
\newcommand{\D}{\mathbb D}
\newcommand{\T}{\mathbb T}
\newcommand{\R}{\mathbb R}
\newcommand{\C}{\mathbb C}
\newcommand{\N}{\mathbb N}
\newcommand{\dd}{\,\mathrm d}
\newcommand{\RH}{\mathrm{RH}}
\DeclareMathOperator{\Rea}{Re}
\DeclareMathOperator{\Arg}{Arg}
\DeclareMathOperator{\Li}{Li}
\lstset{basicstyle=\ttfamily\scriptsize,breaklines=true,columns=fullflexible,
 keepspaces=true,showstringspaces=false,frame=single,rulecolor=\color{black!15},
 backgroundcolor=\color{black!1},keywordstyle=\color{accent}}
\begin{document}
\hypersetup{pageanchor=false}
\begin{titlepage}
\centering
\vspace*{18mm}
{\small\color{accent}\sffamily EXTERNAL APPLICATIONS OF PRIME-CIRCLE METHODS\par}
\vspace{12mm}
{\LARGE\bfseries\color{ink}Toeplitz Counterexamples\par\vspace{3mm}
and Mertens Mean Oscillations\par}
\vspace{7mm}
{\large An exact conjecture refutation and a convex-envelope argument\par}
\vspace{13mm}
{Research prepared for Keio\par}
\vspace{3mm}
{9 November 2026\par}
\vspace{12mm}
\begin{minipage}{0.94\textwidth}
\textbf{Abstract.}
We apply two structures developed in a prime-circle investigation to questions
posed independently in the literature. First, the canonical strongly convex
function of order $0<\alpha<1$ used in Kowalczyk--Lecko--\'{S}miarowska (2021)
has fourth Hermitian Toeplitz determinant
\[
 (1-\alpha^2)^3\left(1-\frac{\alpha^2}{324}\right).
\]
It therefore refutes the lower bound in their Conjecture~1. At
$\alpha=1/2$, an exact rational witness gives $1295/3072$ against the claimed
lower bound $1296/3072$. Toeplitz prediction-error monotonicity propagates
the violation to every matrix order $q\ge4$. Second, a slope-selection lemma
for a concave lower envelope of zero-dependent affine functions, combined
with Zhao's published oscillation and second-moment lemmas, removes the
smooth-boundary assumption in the hypothetical $\Theta=1$ case of his
Mertens-product mean theorem. The prime-power tail is retained exactly, and the other cases are
joined with explicit estimates to complete the original mean-value
characterization of RH.
The constructions, quantifiers, and imported inputs are stated explicitly.
The work does not prove RH. Classical moment positivity and convex-envelope
techniques are credited as such; historical priority of the particular
applications is not certified by the bounded literature search.
\end{minipage}
\vfill
\begin{minipage}{0.94\textwidth}
\centering\small
\fcolorbox{ink!20}{ink!4}{\parbox{0.9\textwidth}{\centering
\textbf{Exact finite certificate}\par\smallskip
$\displaystyle
\det\begin{pmatrix}
1&\tfrac12&\tfrac14&\tfrac7{48}\\
\tfrac12&1&\tfrac12&\tfrac14\\
\tfrac14&\tfrac12&1&\tfrac12\\
\tfrac7{48}&\tfrac14&\tfrac12&1
\end{pmatrix}
=\frac{1295}{3072}<\frac{1296}{3072}
=\left(1-\frac14\right)^3.$}}
\end{minipage}
\end{titlepage}
\hypersetup{pageanchor=true}
\pagenumbering{roman}
\tableofcontents
\clearpage
\pagenumbering{arabic}
\section{Two original targets and the transfer mechanisms}
\label{sec:scope}

The companion manuscript \emph{Prime Reciprocal Circles and the Riemann
Zeta Function: A Unified Energy, Operator, and Sign Theory} developed finite
Toeplitz energies and one-sided arithmetic moments. The present paper uses
those viewpoints on two externally posed questions. All arguments needed
for the transfers are reproduced here, apart from clearly cited classical
theorems and the two analytic estimates stated in Section~\ref{sec:zhao-inputs}.

\begin{center}\small
\begin{tabular}{@{}>{\raggedright\arraybackslash}p{.25\textwidth}>{\raggedright\arraybackslash}p{.31\textwidth}>{\raggedright\arraybackslash}p{.36\textwidth}@{}}
\toprule
Original target & Mechanisms combined & Result and scope\\
\midrule
Kowalczyk et al., Conjecture~1 (2021)
& A disk-to-sector map, exact Taylor coefficients, reflection blocks,
and Toeplitz prediction errors
& The original lower bound is false for every $0<\alpha<1$ and every
$q\ge4$. A rational $4\times4$ witness already proves the refutation.\\[5pt]
Zhao, Section~8, Question~(1) (2025)
& A lower envelope indexed by actual zeros, concavity, local oscillation,
and an upper second-moment bound
& The difficult $\Theta=1$ branch has unbounded mean oscillations in both
signs without the proposed smooth-envelope assumption.\\
\bottomrule
\end{tabular}
\end{center}

The first statement changes neither the analytic class nor the constant in
the proposed inequality. The second concerns the original Mertens product
and its ordinary integral, not a modified sensor. In both cases the
mathematical conclusion is separate from the question of historical
priority. In particular, the general theory of positive Toeplitz matrices
and lower envelopes of zero lines predates this work.

\paragraph{What is not asserted.}
No finite computation is used as evidence for RH. We do not claim a sharp
replacement lower bound for the strongly convex class. The remaining
statements about RH are equivalences or implications proved using the
specified established inputs. The source/status audit at the end records the scope of the proofs,
the imported estimates, and the limits of the historical search.



\section{An exact counterexample to a strongly convex Toeplitz conjecture}
\label{sec:kls-counterexample}

The circle-moment viewpoint has an application to an independently
published conjecture.  Kowalczyk, Lecko and \'Smiarowska
\cite[Conjecture~1, p.~331]{KLS2021} proposed that
\begin{equation}
 (1-\alpha^2)^{q-1}\leq \det T_{q,1}(f)\leq1,
 \qquad 0<\alpha\leq1,\quad q\geq2,
 \label{eq:kls-original}
\end{equation}
for every normalized strongly convex function of order $\alpha$.
The statement is false for every $0<\alpha<1$ and every $q\geq4$.
The first failure has a rational certificate of size four.  The
counterexample below is the very function used in
\cite[Eq.~(21), p.~329]{KLS2021} to attain their proved low-order bounds.
The conclusion does not identify an optimal replacement lower bound.

\subsection{Definitions and an admissible analytic witness}

Write $\mathbb D=\{z\in\mathbb C:|z|<1\}$, and normalize
$f(z)=z+\sum_{n\geq2}a_nz^n$.  Its Hermitian Toeplitz matrix is
\[
 T_{q,1}(f)=
 \begin{pmatrix}
 1&a_2&a_3&\cdots&a_q\\
 \overline{a_2}&1&a_2&\cdots&a_{q-1}\\
 \overline{a_3}&\overline{a_2}&1&\cdots&a_{q-2}\\
 \vdots&\vdots&\vdots&\ddots&\vdots\\
 \overline{a_q}&\overline{a_{q-1}}&\overline{a_{q-2}}&\cdots&1
 \end{pmatrix}.
\]
The class $\mathcal C_\alpha$ in \cite{KLS2021} is defined by
\[
 \left|\arg\left(1+\frac{zf''(z)}{f'(z)}\right)\right|
 \leq\frac{\alpha\pi}{2}\qquad(z\in\mathbb D).
\]
For $\alpha=1$ this is the normalized convex class.

Fix $0<\alpha<1$, take the analytic principal logarithm on the right
half-plane, and define
\begin{align}
 Q_\alpha(z)&=
 \exp\left(\alpha\operatorname{Log}\frac{1+z}{1-z}\right),
 \label{eq:kls-Q}\\
 H_\alpha(z)&=\int_0^z\frac{Q_\alpha(t)-1}{t}\,dt,
 &f_\alpha(z)&=\int_0^z\exp(H_\alpha(w))\,dw.
 \label{eq:kls-witness}
\end{align}
The quotient in the first integral has a removable singularity at zero.
Both integrals are primitives of holomorphic functions on the simply
connected disk and are consequently path independent.

\begin{lemma}[Membership in the original class]
\label{lem:kls-membership}
For every $0<\alpha<1$, the function $f_\alpha$ belongs to
$\mathcal C_\alpha$ and has real Taylor coefficients.  In particular,
it is analytic, convex and univalent throughout $\mathbb D$, with
$f_\alpha(0)=0$ and $f'_\alpha(0)=1$.
\end{lemma}

\begin{proof}
The Cayley map $(1+z)/(1-z)$ takes $\mathbb D$ to the open right
half-plane.  Thus $Q_\alpha$ is analytic, $Q_\alpha(0)=1$, and
$|\arg Q_\alpha(z)|<\alpha\pi/2$.  Differentiating
\eqref{eq:kls-witness} gives $f'_\alpha=\exp H_\alpha\ne0$ and
\[
 1+\frac{zf''_\alpha(z)}{f'_\alpha(z)}=Q_\alpha(z).
\]
The defining sector condition is satisfied, with strict inequality.
The normalization follows directly from $H_\alpha(0)=0$.
Since $\Re Q_\alpha>0$, the standard analytic convexity criterion
also gives convexity and univalence.  Conjugation symmetry of
\eqref{eq:kls-Q} and its primitives proves that the coefficients are real.
\end{proof}

\subsection{A fourth-order defect detected by a Schur complement}

\begin{theorem}[Exact failure of the published lower bound]
\label{thm:kls-four}
For every $0<\alpha<1$,
\begin{equation}
 \det T_{4,1}(f_\alpha)
  =(1-\alpha^2)^3\left(1-\frac{\alpha^2}{324}\right)
  <(1-\alpha^2)^3.
 \label{eq:kls-det4}
\end{equation}
Consequently \cite[Conjecture~1]{KLS2021} is false.
\end{theorem}

\begin{proof}
Expanding the analytic logarithm and exponentials yields
\begin{align*}
 Q_\alpha(z)
  &=1+2\alpha z+2\alpha^2z^2
     +\frac{2\alpha+4\alpha^3}{3}z^3+O(z^4),\\
 H_\alpha(z)
  &=2\alpha z+\alpha^2z^2
     +\frac{2\alpha+4\alpha^3}{9}z^3+O(z^4),\\
 f_\alpha(z)
  &=z+\alpha z^2+\alpha^2z^3
     +\frac{\alpha+17\alpha^3}{18}z^4+O(z^5).
\end{align*}
Set
\[
 a=\alpha,\qquad v=1-a^2,\qquad
 c=\frac{a+17a^3}{18},\qquad
 d=c-a^3=\frac{a(1-a^2)}{18}.
\]
The matrix in question is
\[
 T_{4,1}(f_\alpha)=
 \begin{pmatrix}
 1&a&a^2&c\\
 a&1&a&a^2\\
 a^2&a&1&a\\
 c&a^2&a&1
 \end{pmatrix}.
\]
Its top-left block $A=T_{3,1}(f_\alpha)$ satisfies
\[
 \det A=v^2,
 \qquad
 A^{-1}=\frac1v
 \begin{pmatrix}
 1&-a&0\\
 -a&1+a^2&-a\\
 0&-a&1
 \end{pmatrix}.
\]
For $b=(c,a^2,a)^T$, a direct multiplication gives
\[
 b^TA^{-1}b=a^2+\frac{d^2}{v}.
\]
The Schur-complement determinant formula therefore gives
\[
 \det T_{4,1}(f_\alpha)
  =v^2\left(1-a^2-\frac{d^2}{v}\right)
  =v^3-vd^2
  =v^3\left(1-\frac{a^2}{324}\right).
\]
Since $v>0$ and $a>0$, the strict inequality follows.
\end{proof}

\begin{remark}[Independent two-block verification]
Reflection of coordinates commutes with $T_{4,1}(f_\alpha)$.
Its restrictions to the even and odd reflection subspaces are
\[
 B_+=\begin{pmatrix}1+c&a+a^2\\a+a^2&1+a\end{pmatrix},
 \qquad
 B_-=\begin{pmatrix}1-c&a-a^2\\a-a^2&1-a\end{pmatrix}.
\]
Their determinants are $(1+a)(v+d)$ and $(1-a)(v-d)$.
Their product is $v(v^2-d^2)$, independently recovering
\eqref{eq:kls-det4}.  The defect $d$ enters the two sectors with
opposite signs and hence lowers their product by the strictly positive
quantity $vd^2$.
\end{remark}

\begin{corollary}[A rational certificate]
\label{cor:kls-rational}
For $\alpha=1/2$, an admissible function has
$a_2=1/2$, $a_3=1/4$ and $a_4=7/48$, and
\[
 \det\begin{pmatrix}
 1&1/2&1/4&7/48\\
 1/2&1&1/2&1/4\\
 1/4&1/2&1&1/2\\
 7/48&1/4&1/2&1
 \end{pmatrix}
 =\frac{1295}{3072}
 <\frac{1296}{3072}=\left(1-\frac14\right)^3.
\]
The lower bound fails by exactly $1/3072$.
\end{corollary}

\subsection{Moment positivity and propagation to every higher order}

The propagation step uses only classical facts, which are stated
explicitly to distinguish them from the counterexample calculation.
The Marx--Strohh\"acker theorem gives
\begin{equation}
 \Re\frac{f(z)}z>\frac12\qquad(z\in\mathbb D)
 \label{eq:kls-MS}
\end{equation}
for every normalized convex function $f$; see, for example,
\cite{SugawaWang2015} for the theorem and its stronger generalizations.

\begin{lemma}[Classical Toeplitz positivity and decreasing innovations]
\label{lem:kls-moment}
If $f$ is normalized and convex, every $T_{q,1}(f)$ is positive
semidefinite and
\[
 0\leq D_q(f):=\det T_{q,1}(f)\leq1.
\]
With $D_0=D_1=1$, the positive successive determinant ratios
$D_{q+1}/D_q$ are nonincreasing.  Once a determinant vanishes, all
subsequent determinants vanish.
\end{lemma}

\begin{proof}
By \eqref{eq:kls-MS}, the function $P(z)=2f(z)/z-1$ has positive
real part and $P(0)=1$.  Herglotz representation identifies the
coefficients $a_{n+1}$ with the moments of a positive probability
measure on the unit circle, making $T_{q,1}(f)$ a Gram matrix.
Alternatively, for $0<r<1$, the positive weight
$\Re P(re^{i\theta})\,d\theta/(2\pi)$ gives the matrix with
entries $r^{|j-k|}(T_{q,1})_{jk}$; its limit as $r\uparrow1$ is
positive semidefinite.  Its diagonal entries equal one, so
Hadamard's determinant inequality gives $D_q\leq1$.

Applying the Hadamard--Fischer inequality to the two overlapping
principal submatrices on indices $\{1,\ldots,q\}$ and
$\{2,\ldots,q+1\}$ gives
\[
 D_{q+1}D_{q-1}\leq D_q^2.
\]
For positive determinants this is exactly monotonicity of the
successive ratios.  It is also the usual monotonicity of monic
orthogonal-polynomial norms: multiplying a minimizing polynomial
by the circle coordinate preserves its norm and produces an
admissible competitor at the next degree.

If $T_{q,1}$ is singular, a null vector extended by a zero coordinate
has zero quadratic form in the positive semidefinite $T_{q+1,1}$.
It is therefore a null vector there too.  Induction proves the last
assertion.
\end{proof}

\begin{theorem}[One analytic witness fails at every order $q\geq4$]
\label{thm:kls-allorders}
For every $0<\alpha<1$ and every integer $q\geq4$, the same
function \eqref{eq:kls-witness} satisfies
\begin{equation}
 0\leq\det T_{q,1}(f_\alpha)
 \leq(1-\alpha^2)^{q-1}
     \left(1-\frac{\alpha^2}{324}\right)^{q-3}
 <(1-\alpha^2)^{q-1}.
 \label{eq:kls-allq}
\end{equation}
\end{theorem}

\begin{proof}
Put $v=1-\alpha^2$ and $r=1-\alpha^2/324$.  The coefficient
calculation gives $D_2=v$, $D_3=v^2$ and $D_4=v^3r$.
Lemma~\ref{lem:kls-moment} gives
$D_{k+1}/D_k\leq D_4/D_3=vr$ for $k\geq3$ while these ratios
are positive.  Therefore
\[
 D_q\leq D_3(vr)^{q-3}=v^{q-1}r^{q-3}.
\]
If any determinant vanishes, the same bound remains true for all
higher orders.  Finally $0<r<1$ and $q-3\geq1$, giving the strict
comparison with the proposed lower bound.
\end{proof}


\begin{corollary}[The relative failure grows with the order]
\label{cor:kls-relative}
For every fixed $0<\alpha<1$,
\[
 \lim_{q\to\infty}
 \frac{\det T_{q,1}(f_\alpha)}{(1-\alpha^2)^{q-1}}=0.
\]
\end{corollary}
\begin{proof}
The quotient is nonnegative and, by \eqref{eq:kls-allq}, bounded
above by $(1-\alpha^2/324)^{q-3}$, which tends to zero.
\end{proof}

\begin{corollary}[Classification of the original parameter range]
\label{cor:kls-parameter}
For $0<\alpha\leq1$ and $q\geq2$, the double inequality
\eqref{eq:kls-original} holds for every member of
$\mathcal C_\alpha$ if and only if $q\in\{2,3\}$ or $\alpha=1$.
The upper bound $\det T_{q,1}(f)\leq1$ remains true for every
allowed $(\alpha,q)$.
\end{corollary}

\begin{proof}
The cases $q=2,3$ are the proved bounds in
\cite[Theorems~4 and~5]{KLS2021}.  Every strongly convex function
is convex, so Lemma~\ref{lem:kls-moment} gives the upper bound
and, when $\alpha=1$, the lower bound $0$.  The remaining parameter
range is refuted by Theorem~\ref{thm:kls-allorders}. The upper
endpoint is attained by $f(z)=z$; when $\alpha=1$, the lower
endpoint is attained by $f(z)=z/(1-z)$.
\end{proof}

\subsection{What this application does and does not establish}

The named statement has been kept unchanged: its original function
class, parameter range, normalization and constants all enter the
counterexample.  An exact analytic function supplies the coefficients,
so membership is not inferred from a truncated Taylor polynomial.
The computation is rational and symbolic; no finite-precision sign
test or unproved asymptotic premise is used.  The all-order conclusion
closes the integer quantifier as well.

The mechanism is a coefficient defect relative to a geometric
moment sequence.  The low-order values $a_2=\alpha$ and
$a_3=\alpha^2$ agree with that sequence, but
$a_4-\alpha^3=\alpha(1-\alpha^2)/18\ne0$.  A Schur complement
turns this defect into a strictly negative correction to the
determinant, and classical prediction-error monotonicity propagates
the correction.  This is an external disproof, not an implication
for the Riemann hypothesis.

Historical priority remains a separate question.  The original
conjecture is verified in the primary 2021 publication, and targeted
searches of subsequent publications and the exact determinant
expression did not identify an earlier matching refutation.  Such
a search does not establish absence from all prior literature.
No claim of first discovery, a historically new general principle,
or complete satisfaction of a no-prior-duplication requirement is
made here.  In particular, the universal convex-class estimate
$0\leq\det T_{q,1}(f)\leq1$ is explicitly presented as a classical
corollary, even though related finite-order conjectures appear in
\cite{Cudna2020,LeckoSimSmiarowska2020}.



\section{The mean error in Mertens' product: a concrete application}
\label{sec:zhao-application}

Put
\[
 B_{\mathbb P}=\gamma_E-
 \sum_p\sum_{k\geq2}\frac{1}{kp^k},\qquad
 E(x)=\sum_{p\leq x}\frac1p-\log\log x-B_{\mathbb P},
\]
and define
\[
 \Delta_1(X)=\int_2^X E(x)\,dx,\qquad
 \Delta_2(X)=\int_2^X E(x)^2\,dx,
\]
\[
 E_3(x)=\frac{1}{\log x}
          \prod_{p\leq x}(1-p^{-1})^{-1}-e^{\gamma_E},
 \qquad J(X)=\int_2^X E_3(x)\,dx.
\]
Here and below a nontrivial zero is denoted by
$\rho=\beta+i\gamma$, and
\[
 \Theta=\sup\{\beta:\zeta(\beta+i\gamma)=0,
                              \ 0<\beta<1\}.
\]

The precise external question is Question~(1) in
\cite[Section~8]{Zhao2025}: can the additional regularity assumption
in the $\Theta=1$ case of that paper's Theorem~2 be removed?
The original question concerns the sign of $J(X)$, with the normalization
above. We retain that function and that normalization.

\begin{theorem}[Removal of the auxiliary zero-boundary assumption]
\label{thm:zhao-removal}
If $\Theta=1$, then
\[
 \limsup_{X\to\infty}J(X)=+\infty,
 \qquad
 \liminf_{X\to\infty}J(X)=-\infty.
\]
No differentiability, regular variation, or sharpness assumption on an
auxiliary description of the zero-free region is required.
More precisely, under the same hypothesis,
\begin{equation}
 \label{eq:zhao-moment-domination}
 \limsup_{X\to\infty}
       \frac{\Delta_1(X)}{1+\Delta_2(X)}=+\infty,
 \qquad
 \liminf_{X\to\infty}
       \frac{\Delta_1(X)}{1+\Delta_2(X)}=-\infty.
\end{equation}
\end{theorem}

The proof below is a deduction from the two published estimates stated
in Section~\ref{sec:zhao-inputs}. It is not a proof of the Riemann
hypothesis. Its application is to the previously stated question of
removing an auxiliary hypothesis from a mean-value characterization.

The zero envelope used here is established prior art. In particular,
Pintz~\cite{Pintz2017} studies
$\inf_\rho\bigl((1-\beta)\log x+\log|\rho|\bigr)$;
R\'ev\'esz~\cite{Revesz2025} discusses convexification and the resulting
concave envelopes. Broucke~\cite[Introduction and Section~2]{Broucke2026}
reviews these connections. We make no claim to have introduced this
envelope or the associated convex-analysis principle. The step used
for the present application is a selection of supporting zeros at
which the first-moment oscillation exceeds the entire second moment.
Priority of this particular application is a separate bibliographic
question; the proof does not establish a historical firstness claim.

\subsection{The two published analytic inputs}
\label{sec:zhao-inputs}

We use the following statements with their quantifiers unchanged.

\begin{lemma}[Zero-specific oscillation; Zhao, Lemma~10]
\label{lem:zhao-published-oscillation}
For every $\varepsilon>0$ there are effectively computable constants
$C_1(\varepsilon),C_2(\varepsilon)$ with the following property.
If $\rho_0=\beta_0+i\gamma_0$ is a zero of $\zeta$ with
\[
 \beta_0>\tfrac12+\varepsilon,
 \qquad \gamma_0>C_1(\varepsilon),
\]
then for every $X>\gamma_0^{C_2(\varepsilon)}$ there exist
$x_+,x_-\in[X,X^{1+\varepsilon}]$ such that
\begin{equation}
 \label{eq:zhao-published-omega}
 \Delta_1(x_+)>
     \frac{x_+^{\beta_0}}{\gamma_0^{2+\varepsilon}\log x_+},
 \qquad
 \Delta_1(x_-)<-
     \frac{x_-^{\beta_0}}{\gamma_0^{2+\varepsilon}\log x_-}.
\end{equation}
\end{lemma}

\begin{lemma}[Second moment; Zhao, Lemma~11]
\label{lem:zhao-published-moment}
Suppose $\eta:[0,\infty)\to(0,1/2]$ is continuous and
nonincreasing, and that $\zeta(s)$ has no zeros in
$\operatorname{Re}s>1-\eta(|\operatorname{Im}s|)$.
Define
\[
 \omega_\eta(x)=\inf_{t\geq1}
                  \{\eta(t)\log x+\log t\}.
\]
For every fixed $0<\varepsilon<1$,
\begin{equation}
 \label{eq:zhao-published-second}
 \Delta_2(x)\ll_\varepsilon
   \frac{x\exp\{-2(1-\varepsilon)\omega_\eta(x)\}}
        {(\log x)^2}.
\end{equation}
The version published in \cite{Zhao2025} also allows finitely many
exceptional zeros when $\Theta=1$; that extension is not needed here.
\end{lemma}

These are imported results, not claims re-proved in this section.
Zhao derives the first by adapting Pintz's localized oscillation method
\cite{Pintz1980I}, and the second from an explicit-formula second-moment
bound and Pintz's zero-free-region estimate \cite{Pintz1980II}.
In particular, the proof of the first lemma in \cite{Zhao2025} is a
sketch; we have not independently reconstructed every step of that
Pintz-method adaptation. Thus the logical dependency of
Theorem~\ref{thm:zhao-removal} includes the published statements of
these two lemmas. Continuity and nonincrease, with no derivative
condition, are the hypotheses in the second lemma.

\subsection{A supporting-line selection lemma}

\begin{lemma}
\label{lem:zhao-selection}
Let $\{(a_j,v_j):j\in\mathcal I\}$ be a nonempty family satisfying
\[
 a_j>0,\quad v_j>0,\quad
 \inf_j a_j=0,\quad a_jv_j\geq c>0,
\]
and suppose that only finitely many indices have $v_j\leq V$ for
any fixed $V$. Fix $a_0>0$ and set
\[
 \Omega(u)=\min\left\{a_0u,\inf_{j\in\mathcal I}(a_ju+v_j)\right\},
 \qquad u>0.
\]
For each $0<\kappa<1$ there are arbitrarily large $u$ and an index
$j$ such that
\begin{equation}
 \label{eq:zhao-selected-zero}
 \Omega(u)=a_ju+v_j,\qquad a_ju\geq\kappa v_j.
\end{equation}
Along any such sequence with $u\to\infty$,
\[
 a_j\to0,\qquad v_j\to\infty,\qquad v_j=o(u).
\]
\end{lemma}

\begin{proof}
For fixed $u$, the auxiliary line bounds the infimum by $a_0u$;
therefore only indices with $v_j\leq a_0u$ can minimize it. There
are finitely many of them, so the minimum is attained.
As the infimum of positive affine functions with positive slopes,
$\Omega$ is positive, nondecreasing, and concave, hence locally
absolutely continuous on $(0,\infty)$. The arithmetic--geometric
mean inequality gives
\[
 a_ju+v_j\geq2\sqrt{a_jv_ju}\geq2\sqrt{cu}.
\]
The auxiliary line also exceeds a positive constant times $\sqrt u$
for all sufficiently large $u$. Consequently
\begin{equation}
 \label{eq:zhao-envelope-lower}
 \Omega(u)\gg\sqrt u\qquad(u\to\infty).
\end{equation}
On the other hand, $\Omega(u)/u\leq a_j+v_j/u$ for every fixed
$j$, and $\inf_j a_j=0$. Thus
\begin{equation}
 \label{eq:zhao-envelope-sublinear}
 \Omega(u)=o(u).
\end{equation}
In particular the auxiliary line is not a minimizer for all
sufficiently large $u$.

Suppose the asserted selection fails. Then for all sufficiently
large $u$, every minimizing actual line satisfies $a_ju<\kappa v_j$.
At a point of differentiability of $\Omega$, the slope of every
minimizing line equals $\Omega'(u)$. Hence, almost everywhere on a
terminal interval,
\[
 \frac{u\Omega'(u)}{\Omega(u)}
 =\frac{a_ju}{a_ju+v_j}
 <\frac{\kappa}{1+\kappa}.
\]
Integrating the logarithmic derivative gives
\[
 \Omega(u)\ll u^{\kappa/(1+\kappa)}.
\]
Since $\kappa/(1+\kappa)<1/2$, this contradicts
\eqref{eq:zhao-envelope-lower}.

For a selected pair, $a_j\leq\Omega(u)/u\to0$ and
$v_j\leq\Omega(u)=o(u)$. If $v_j$ stayed bounded along an
infinite subsequence, local finiteness would put those indices in
a fixed finite set with strictly positive minimum $a_j$, a
contradiction. This proves the remaining conclusions.
\end{proof}

\subsection{One fixed continuous zero-free envelope}

Assume now that $\Theta=1$. Let $\mathcal Z$ be the nontrivial
zeros with $\gamma>0$ and $\beta>3/4$, counted with multiplicity,
and put
\[
 a_\rho=1-\beta,\qquad v_\rho=\log\gamma,
 \qquad a_0=\tfrac14.
\]
The familiar zero-free region
$\beta\leq1-c_0/\log(\gamma+2)$ gives
$a_\rho v_\rho\geq c>0$, after decreasing $c$ if necessary.
Here the ordinates are bounded away from $1$, and each bounded
height contains only finitely many zeros. The hypothesis $\Theta=1$
gives $\inf_{\rho\in\mathcal Z}a_\rho=0$.
Thus Lemma~\ref{lem:zhao-selection} applies to
\begin{equation}
 \label{eq:zhao-discrete-envelope}
 \Omega(u)=\min\left\{\frac u4,
           \inf_{\rho\in\mathcal Z}
                  ((1-\beta)u+\log\gamma)\right\}.
\end{equation}

For $v\geq0$ define the running boundary
\[
 b(v)=\min\left\{\tfrac14,
       \inf_{\substack{\rho\in\mathcal Z\\v_\rho\leq v}}a_\rho
             \right\},
\]
where the infimum over an empty set is ignored. This is a positive
nonincreasing function, piecewise constant on compact intervals.
Directly from the definition,
\begin{equation}
 \label{eq:zhao-running-envelope}
 \inf_{v\geq0}\{b(v)u+v\}=\Omega(u).
\end{equation}
Indeed, a record value $a_\rho$ begins at $v=v_\rho$; waiting
longer at the same value increases $b(v)u+v$. The initial value
$1/4$ gives the auxiliary line at $v=0$.

Define, once and for all,
\begin{equation}
 \label{eq:zhao-smoothed-eta}
 \eta(t)=\int_0^1 b(\log t+s)\,ds\quad(t\geq1),
 \qquad \eta(t)=\eta(1)\quad(0\leq t\leq1).
\end{equation}
This function is positive, continuous, and nonincreasing, and
\begin{equation}
 \label{eq:zhao-eta-sandwich}
 b(\log t+1)\leq\eta(t)\leq b(\log t)
 \qquad(t\geq1).
\end{equation}
At every zero in $\mathcal Z$,
$\eta(\gamma)\leq b(\log\gamma)\leq1-\beta$.
For the remaining zeros $\beta\leq3/4$, the same inequality
follows from $\eta\leq1/4$. Conjugation handles negative ordinates.
It follows that the open region
$\sigma>1-\eta(|t|)$ is zero free.

This construction uses one fixed $\eta$, independent of the selected
zeros and of $u$. In particular the implied constant in
Lemma~\ref{lem:zhao-published-moment} does not change with the
subsequence to be chosen. Equations
\eqref{eq:zhao-running-envelope}--\eqref{eq:zhao-eta-sandwich} give
\begin{equation}
 \label{eq:zhao-omega-comparison}
 \Omega(u)-1\leq\omega_\eta(e^u)\leq\Omega(u).
\end{equation}
For the lower bound, write $v=\log t$ and observe that
\[
 \eta(e^v)u+v\geq b(v+1)u+v
  =b(v+1)u+(v+1)-1\geq\Omega(u)-1.
\]
The upper bound follows by taking infima in
$\eta(e^v)\leq b(v)$.

\subsection{First-moment oscillations dominate the second moment}

Fix the numerical choices
\[
 \kappa=\frac12,\qquad\varepsilon=\frac1{20}.
\]
By Lemma~\ref{lem:zhao-selection}, there are $u_j\to\infty$ and
zeros $\rho_j=\beta_j+i\gamma_j$ for which, with
$a_j=1-\beta_j$ and $v_j=\log\gamma_j$,
\begin{equation}
 \label{eq:zhao-chosen-parameters}
 \Omega(u_j)=a_ju_j+v_j,\qquad
 a_ju_j\geq\tfrac12v_j,\qquad
 a_j\to0,\quad v_j\to\infty,\quad v_j=o(u_j).
\end{equation}
In particular, for all sufficiently large $j$,
\[
 \beta_j>\tfrac12+\varepsilon,\qquad
 \gamma_j>C_1(\varepsilon),\qquad
 e^{u_j}>\gamma_j^{C_2(\varepsilon)}.
\]
Lemma~\ref{lem:zhao-published-oscillation} supplies points
$x_{j,+},x_{j,-}\in[e^{u_j},e^{(1+\varepsilon)u_j}]$ of the
respective signs. Write $w_{j,\pm}=\log x_{j,\pm}$.
For either sign, abbreviate $w=w_{j,\pm}$ and
\[
 A_{j,\pm}=|\Delta_1(e^w)|.
\]
Then
\begin{equation}
 \label{eq:zhao-amplitude-lower}
 A_{j,\pm}>
 \frac{\exp\{w-a_jw-(2+\varepsilon)v_j\}}{w}.
\end{equation}
The fixed-envelope estimate
\eqref{eq:zhao-published-second}, the comparison
\eqref{eq:zhao-omega-comparison}, and monotonicity of $\Omega$
give, with a fixed constant $C>0$,
\begin{equation}
 \label{eq:zhao-moment-upper}
 \Delta_2(e^w)\leq
 C\frac{\exp\{w-2(1-\varepsilon)(a_ju_j+v_j)\}}{w^2}.
\end{equation}
Divide \eqref{eq:zhao-amplitude-lower} by
\eqref{eq:zhao-moment-upper}. Since
$w\leq(1+\varepsilon)u_j$,
\begin{align}
 \log\frac{A_{j,\pm}}{\Delta_2(e^w)}
 &\geq
  2(1-\varepsilon)a_ju_j-a_jw
  -3\varepsilon v_j+\log w-O(1)\notag\\
 &\geq
  (1-3\varepsilon)a_ju_j
  -3\varepsilon v_j+\log w-O(1)\notag\\
 &\geq
  \left(\tfrac12(1-3\varepsilon)-3\varepsilon\right)v_j
  +\log w-O(1)\notag\\
 &=\frac{11}{40}v_j+\log w-O(1)
   \longrightarrow+\infty.
 \label{eq:zhao-final-ratio}
\end{align}
If $\Delta_2(e^w)=0$, the required domination is immediate;
in this application it is positive for every $e^w>2$.
Furthermore, since $a_j\to0$ and $v_j=o(u_j)$,
\eqref{eq:zhao-amplitude-lower} implies
\begin{equation}
 \label{eq:zhao-amplitude-divergence}
 A_{j,\pm}\geq\frac{\exp\{(1-o(1))u_j\}}
                         {(1+\varepsilon)u_j}
 \longrightarrow\infty.
\end{equation}
Equations \eqref{eq:zhao-final-ratio} and
\eqref{eq:zhao-amplitude-divergence} prove
\eqref{eq:zhao-moment-domination}.

The distinction from merely changing variables is visible in
\eqref{eq:zhao-final-ratio}: it is an asymptotically divergent ratio
between the signed contribution and the complete nonnegative
quadratic error. No cancellation conjecture or finite-height zero
verification enters this comparison.

\subsection{The exact product identity and both signs}

The conversion from $E$ to $E_3$ must retain the prime-power tail
\begin{equation}
 \label{eq:zhao-product-tail}
 R(x)=\sum_{p>x}\sum_{k\geq2}\frac{1}{kp^k}\geq0.
\end{equation}
The exact Euler-product identity is
\begin{equation}
 \label{eq:zhao-exact-product}
 e^{-\gamma_E}E_3(x)=\exp\{E(x)-R(x)\}-1.
\end{equation}
Indeed, expanding the logarithm of the finite product and using the
definition of $B_{\mathbb P}$ gives
\[
 \log\left(\frac{e^{-\gamma_E}}{\log x}
             \prod_{p\leq x}(1-p^{-1})^{-1}\right)=E(x)-R(x).
\]
The tail is small by an elementary bound:
\begin{equation}
 \label{eq:zhao-tail-bound}
 0\leq R(x)\leq
 \frac12\sum_{n>x}\frac1{n(n-1)}
 =\frac{1}{2\lfloor x\rfloor}\leq\frac1x
 \qquad(x\geq2).
\end{equation}
Here $1/k\leq1/2$ for $k\geq2$ was applied before summing
the geometric series, and replacing primes by all integers only
increases the sum. In particular
\begin{equation}
 \label{eq:zhao-tail-integral}
 \int_2^X R(x)\,dx\leq\log(X/2).
\end{equation}

By Mertens' theorem $E(x)\to0$; local boundedness then gives
$|E(x)|\leq M$ for all $x\geq2$ and some finite $M$.
Taylor's theorem implies
\[
 e^E-1\leq E+C E^2\qquad(|E|\leq M)
\]
with one fixed $C$ (for example $C=e^M/2$).
Using $R\geq0$ for the upper bound and $e^y-1\geq y$ for
the lower bound gives the valid two-sided inequalities
\begin{equation}
 \label{eq:zhao-exact-mean-sandwich}
 \Delta_1(X)-\log(X/2)
 \leq e^{-\gamma_E}J(X)
 \leq\Delta_1(X)+C\Delta_2(X).
\end{equation}
At $X=x_{j,-}$, the right-hand side tends to $-\infty$ by
\eqref{eq:zhao-final-ratio} and
\eqref{eq:zhao-amplitude-divergence}. At $X=x_{j,+}$, the
left-hand side tends to $+\infty$: the exponential lower bound
in \eqref{eq:zhao-amplitude-divergence} exceeds
$\log x_{j,+}=O(u_j)$. This proves
Theorem~\ref{thm:zhao-removal}. Since $J$ is continuous, the two
unbounded sequences of opposite signs also imply infinitely many
sign changes.

\begin{remark}[Exactness of the nonlinear conversion]
The displayed exponential conversion in \cite[Section~4]{Zhao2025}
does not include the term $R(x)$. With the definitions of that paper,
\eqref{eq:zhao-exact-product} is the exact identity. The negative
tail helps the upper bound needed at the negative points; its
integral is negligible at the positive points by
\eqref{eq:zhao-tail-integral}. Thus the argument above does not
rely on dropping this term or on an incorrect exact equality.
\end{remark}


\subsection{The full characterization with the tail retained}

\begin{corollary}[The complete mean-value characterization, with the tail retained]
\label{cor:zhao-complete-characterization}
With the original Mertens-product error $E_3$ and the original integral $J$,
\[
 \mathrm{RH}
 \quad\Longleftrightarrow\quad
 J(X)>0\quad\hbox{for every }X>2.
\]
If RH is false, then
\[
 \limsup_{X\to\infty}J(X)=+\infty,
 \qquad
 \liminf_{X\to\infty}J(X)=-\infty.
\]
The source estimates used for the $\Theta<1$ cases are specified explicitly
in the proof; the exponential identity with an omitted prime-power tail
is not used.
\end{corollary}

\begin{proof}
We use the exact sandwich
\eqref{eq:zhao-exact-mean-sandwich}, already established with the
prime-power tail retained.

\emph{RH implies positivity for every $X>2$.}
Rosser and Schoenfeld's Theorem~23, specifically the lower inequality in
\cite[Eq.~(4.10), p.~73]{RosserSchoenfeld1962}, states
\[
 e^{\gamma_E}\log x<\prod_{p\leq x}\frac{p}{p-1}
 \qquad(0<x\leq10^8).
\]
Thus $E_3(x)>0$ for $2\leq x\leq10^8$, and hence $J(X)>0$ for
$2<X\leq10^8$.

For the remaining range we use the quantitative estimate proved at the
end of \cite[Section~2, following Eq.~(19)]{Zhao2025}. Under RH,
\begin{equation}
 \label{eq:zhao-rh-imported-tail-bound}
 \int_X^\infty\frac{\pi(t)-\operatorname{li}(t)}{t^2}\,dt
 <-\frac{1.73}{\sqrt X\log X}
 \qquad(X\geq10^8).
\end{equation}
This estimate belongs to that paper's first-moment proof and does not
use its later exponential conversion. The exact partial-summation identity
\cite[Eq.~(18)]{Zhao2025} is
\[
 \Delta_1(X)
 =-X\int_X^\infty\frac{\pi(t)-\operatorname{li}(t)}{t^2}\,dt+C_0,
 \qquad
 C_0=2B_{\mathbb P}+2\log\log2-\operatorname{li}(2).
\]
In particular,
\begin{equation}
 \label{eq:zhao-rh-first-moment-explicit}
 \Delta_1(X)>1.73\frac{\sqrt X}{\log X}-2
 \qquad(X\geq10^8).
\end{equation}
For completeness, the rough constant bound $C_0>-2$ needs no numerical
prime computation. Indeed,
\[
 \sum_p\sum_{k\geq2}\frac1{kp^k}
 \leq\sum_{n\geq2}\left[-\log(1-n^{-1})-n^{-1}\right]
 =1-\gamma_E,
\]
so $B_{\mathbb P}\geq2\gamma_E-1$. If $t=\log2$, then
\[
 \operatorname{li}(2)=\gamma_E+\log t+
       \sum_{k\geq1}\frac{t^k}{k\,k!},
 \qquad
 \sum_{k\geq1}\frac{t^k}{k\,k!}<e^t-1=1.
\]
The elementary inequalities $\gamma_E>1/2$ and $\log t>-1/2$ give
\[
 C_0\geq3\gamma_E-2+\log t-
              \sum_{k\geq1}\frac{t^k}{k\,k!}>-2.
\]

For $X\geq10^8$, the function $\sqrt X/\log X$ is increasing and
exceeds $500$. Thus \eqref{eq:zhao-rh-first-moment-explicit} gives
$\Delta_1(X)>\sqrt X/\log X$. Also
\[
 \frac{\sqrt X}{\log X}>\log X\qquad(X\geq10^8).
\]
To verify the last inequality, write $u=\log X$: $e^{u/2}/u^2$ is
increasing for $u>4$ and exceeds $1$ at $X=10^8$.
Using the lower inequality in \eqref{eq:zhao-exact-mean-sandwich}, we obtain
\[
 e^{-\gamma_E}J(X)
 >\log X-\log(X/2)=\log2>0
 \qquad(X\geq10^8).
\]
This proves the entire RH direction while retaining $R$.

\emph{The case $1/2<\Theta<1$.}
The two inputs are
\cite[Eq.~(20)]{Zhao2025}, which gives for every $\varepsilon>0$
\[
 \Delta_1(X)=\Omega_\pm(X^{\Theta-\varepsilon}),
\]
and \cite[Lemma~9]{Zhao2025}, which gives
\[
 \Delta_2(X)=O_\Theta\!\left(
       \frac{X^{2\Theta-1}}{(\log X)^2}\right).
\]
Choose $0<\varepsilon<1-\Theta$. Then
$\Theta-\varepsilon>2\Theta-1>0$. Along the positive sequence
provided by the first input, the lower inequality in
\eqref{eq:zhao-exact-mean-sandwich} tends to $+\infty$, because
$\log X=o(X^{\Theta-\varepsilon})$. Along its negative sequence,
the upper inequality tends to $-\infty$, because
$\Delta_2(X)=o(X^{\Theta-\varepsilon})$. Thus the discarded tail in
the original exponential calculation does not change either sign in
this case.

\emph{The case $\Theta=1$.}
Theorem~\ref{thm:zhao-removal} gives the two unbounded signs without
the additional regularity assumption. By the functional equation,
failure of RH is equivalent to $\Theta>1/2$. The preceding two cases
therefore establish the converse and the stated oscillation assertion.
\end{proof}

\begin{remark}[Scope of the imported estimates]
The proof above independently repairs the nonlinear passage to the
Mertens product. It imports Rosser--Schoenfeld's finite-range product
inequality, Zhao's explicit first-moment bound under RH, and Zhao's
first- and second-moment estimates for $1/2<\Theta<1$.
The $\Theta=1$ argument retains its separately stated dependency on the
published localized oscillation and zero-free-region second-moment
lemmas. None of these dependencies is replaced by the incorrect
identity obtained by setting $R$ equal to zero.
\end{remark}

\section{Scope, reproducibility, and historical status}
\label{sec:audit}

\subsection{What the mathematical arguments close}

The Toeplitz conclusion is an unconditional refutation of a universally
quantified assertion in an identifiable published source. A counterexample
is a complete negative answer to that assertion. The analytic witness is
defined on the entire disk, not only by a finite coefficient list. The
fourth-order determinant can be checked independently of every result in
the number-theoretic part of this paper. The higher-order argument also
classifies the whole original parameter range, while leaving the problem
of a replacement sharp lower bound open here.

The Mertens argument addresses the original integrated product error.
Its new deduction from the specified published estimates is the existence
of supporting zeros for which the signed first moment dominates the full
second moment. The envelope is fixed before the subsequence is chosen, so
the implied constants do not depend on a varying auxiliary function. All
selected points tend to infinity, and the prime-power tail is bounded
before drawing conclusions about the product.

\begin{center}\small
\begin{tabular}{@{}>{\raggedright\arraybackslash}p{.23\textwidth}>{\raggedright\arraybackslash}p{.69\textwidth}@{}}
\toprule
Verification question & Finding\\
\midrule
Externally posed target & KLS Conjecture~1 and Zhao Section~8,
Question~(1), both identified in primary publications.\\[4pt]
Original problem retained & The analytic class, matrix normalization,
parameters and constants are unchanged in the first target. The second
uses the original Mertens product and ordinary integral.\\[4pt]
Complete quantifiers & The witness fails for all $q\ge4$ and all
$0<\alpha<1$. In the $\Theta=1$ branch, both signed sequences are
unbounded and no additional regularity hypothesis is assumed.\\[4pt]
Exact versus numerical & The Toeplitz calculation is a polynomial identity
over $\mathbb Q$. The analytic application uses estimates and a subsequence
argument, not numerical zero data.\\[4pt]
Imported theory & Classical convexity, Herglotz positivity and
Hadamard--Fischer; for the arithmetic application, Zhao's stated local
oscillation and second-moment lemmas and the classical zero-free region.\\[4pt]
Novelty and open status & The sources explicitly pose the targets.
No matching resolution was identified in the bounded subsequent-literature
search. That does not certify continued open status at the time of this
note or historical priority of these arguments.\\
\bottomrule
\end{tabular}
\end{center}

\subsection{The relation to the originating prime-circle manuscript}

The transferable objects are positive circle moments and signed arithmetic
moments. In the first application, circle moments supply a sequence of
positive Toeplitz matrices. A coefficient discrepancy that is invisible
to the first two determinants is detected by a fourth-order Schur
complement, then persists through decreasing prediction errors. In the
second application, a convex-analytic description of zero locations is
combined with first- and second-moment estimates; the useful conclusion is
the divergent ratio in \eqref{eq:zhao-final-ratio}.

These are specific applications of established mathematical structures.
The Cayley map, Schur complements, prediction errors and lower envelopes
are not introduced as historically new principles. Their names do not
establish a transfer on their own: the transfer is the proved class
membership and negative determinant correction in the first case, and the
uniform estimates and supporting-line selection in the second.

\subsection{Literature checks and remaining limits}

The literature check used the original authors, titles, conjecture wording,
DOIs, exact determinant constants, and later work on strongly convex
Hermitian Toeplitz determinants. An indexed citation list for KLS included
eleven later records dated 2022--2026; the inspected nearby papers concerned
different function classes or low-order determinants. Searches for Zhao's
question included its assumption number, integrated Mertens errors,
first/second moments and zero envelopes. No matching refutation or
assumption removal was found in the inspected material.

This is a reproducible description of a bounded search, not a proof of
absence from the literature. In particular, this paper does not assert
that all historical novelty requirements have been certified. The general
envelope method is positively identified as prior art in
\cite{Pintz2017,Revesz2025,Broucke2026}. The original strongly convex
witness is credited to \cite{KLS2021}. The claims should be judged on the
explicit proofs and these stated dependencies, not on an unsupported
claim of first discovery.

\appendix
\section{An exact arithmetic certificate}
\label{app:certificate}

The following standalone Python program uses only the standard library
and rational arithmetic. It verifies the complete polynomial identity by
the permutation formula, separately from the Schur-complement proof.
It then checks the rational instance and its exact gap. It is included
in the single \texttt{.tex} source so that this note has no external code
or figure dependencies. This certificate verifies the arithmetic;
Lemma~\ref{lem:kls-membership} supplies the independent analytic
membership proof.

\begin{lstlisting}[language=Python]
"""Exact rational certificate for the q=4 KLS 2021 counterexample.

No floating-point arithmetic and no external dependencies are used.
"""
from fractions import Fraction as F
from itertools import permutations

def det(M):
    n = len(M)
    out = F(0)
    for p in permutations(range(n)):
        sign = (-1) ** sum(p[i] > p[j] for i in range(n) for j in range(i + 1, n))
        term = F(sign)
        for i in range(n):
            term *= M[i][p[i]]
        out += term
    return out

def add(a, b):
    c = [F(0)] * max(len(a), len(b))
    for i, x in enumerate(a): c[i] += x
    for i, x in enumerate(b): c[i] += x
    while len(c) > 1 and c[-1] == 0: c.pop()
    return c

def mul(a, b):
    c = [F(0)] * (len(a) + len(b) - 1)
    for i, x in enumerate(a):
        for j, y in enumerate(b): c[i+j] += x*y
    while len(c) > 1 and c[-1] == 0: c.pop()
    return c

def pdet(M):
    n = len(M)
    out = [F(0)]
    for p in permutations(range(n)):
        sign = (-1) ** sum(p[i] > p[j] for i in range(n) for j in range(i+1,n))
        term = [F(sign)]
        for i in range(n): term = mul(term, M[i][p[i]])
        out = add(out, term)
    return out

one, a, a2 = [F(1)], [F(0),F(1)], [F(0),F(0),F(1)]
c = [F(0), F(1,18), F(0), F(17,18)]
M = [[one,a,a2,c],[a,one,a,a2],[a2,a,one,a],[c,a2,a,one]]
actual = pdet(M)
v = [F(1),F(0),F(-1)]
claimed = mul(mul(mul(v,v),v),[F(1),F(0),F(-1,324)])
assert actual == claimed
print("Polynomial identity verified exactly:")
print("det T4 = (1 - alpha^2)^3 (1 - alpha^2/324)")
print("Coefficients [degree: coefficient]:", {i:str(x) for i,x in enumerate(actual) if x})

a = F(1,2)
c = (a + 17*a**3)/18
M = [[F(1),a,a*a,c],[a,F(1),a,a*a],[a*a,a,F(1),a],[c,a*a,a,F(1)]]
D = det(M)
B = (1-a*a)**3
assert c == F(7,48)
assert D == F(1295,3072)
assert B == F(1296,3072)
assert B-D == F(1,3072)
print("alpha=1/2 coefficients:", [str(x) for x in M[0]])
print("determinant:", D)
print("conjectured lower bound:", B)
print("exact lower-bound violation:", B-D)
\end{lstlisting}


\section{Reproducibility notes}

The document is self-contained as a LaTeX source and may be compiled with
\texttt{pdflatex} (twice for cross-references) or \texttt{latexmk -pdf}.
There are no external images, bibliography databases or source-file inputs.
All mathematical displays are vector text. The finite certificate has been
run successfully with exact rational arithmetic; the main determinant
identity was also checked by the reflection-block factorization.


\phantomsection
\addcontentsline{toc}{section}{References}
\begin{thebibliography}{99}

\bibitem{KLS2021}
B.~Kowalczyk, A.~Lecko and B.~\'Smiarowska,
\emph{Sharp inequalities for Hermitian Toeplitz determinants for
strongly starlike and strongly convex functions},
Journal of Mathematical Inequalities \textbf{15} (2021), 323--332.
\href{https://doi.org/10.7153/jmi-2021-15-24}{doi:10.7153/jmi-2021-15-24}.
\href{https://files.ele-math.com/articles/jmi-15-24.pdf}{Primary full text}.

\bibitem{SugawaWang2015}
T.~Sugawa and L.-M.~Wang,
\emph{Notes on convex functions of order $\alpha$},
arXiv:1502.05127 (2015).
\href{https://arxiv.org/abs/1502.05127}{arXiv:1502.05127}.

\bibitem{Cudna2020}
K.~Cudna, O.~S.~Kwon, A.~Lecko, Y.~J.~Sim and B.~\'Smiarowska,
\emph{The second and third-order Hermitian Toeplitz determinants
for starlike and convex functions of order $\alpha$},
Bolet\'in de la Sociedad Matem\'atica Mexicana \textbf{26}
(2020), 361--375.
\href{https://doi.org/10.1007/s40590-019-00271-1}{doi:10.1007/s40590-019-00271-1}.

\bibitem{LeckoSimSmiarowska2020}
A.~Lecko, Y.~J.~Sim and B.~\'Smiarowska,
\emph{The fourth-order Hermitian Toeplitz determinant for convex
functions}, Analysis and Mathematical Physics \textbf{10}
(2020), article 39.
\href{https://doi.org/10.1007/s13324-020-00382-3}{doi:10.1007/s13324-020-00382-3}.

\bibitem{Zhao2025}
T.~Zhao,
\emph{On the mean values of the error terms in Mertens' theorems},
Research in Number Theory \textbf{11} (2025), article 62.
\href{https://doi.org/10.1007/s40993-025-00640-y}{doi:10.1007/s40993-025-00640-y}.
\href{https://arxiv.org/html/2411.18903v2}{arXiv:2411.18903v2}.

\bibitem{RosserSchoenfeld1962}
J.~B.~Rosser and L.~Schoenfeld,
\emph{Approximate formulas for some functions of prime numbers},
Illinois Journal of Mathematics \textbf{6} (1962), 64--94.
\href{https://doi.org/10.1215/ijm/1255631807}{doi:10.1215/ijm/1255631807}.
\href{https://denisevellachemla.eu/Rosser-Schoenfeld-1962.pdf}{Primary scan}.

\bibitem{Pintz1980I}
J.~Pintz,
\emph{On the remainder term of the prime number formula I.
On a problem of Littlewood},
Acta Arithmetica \textbf{36} (1980), 341--365.
\href{https://doi.org/10.4064/aa-36-4-341-365}{doi:10.4064/aa-36-4-341-365}.

\bibitem{Pintz1980II}
J.~Pintz,
\emph{On the remainder term of the prime number formula II.
On a theorem of Ingham},
Acta Arithmetica \textbf{37} (1980), 209--220.
\href{https://doi.org/10.4064/aa-37-1-209-220}{doi:10.4064/aa-37-1-209-220}.

\bibitem{Pintz2017}
J.~Pintz,
\emph{Distribution of zeta zeros and the oscillation of the error
term of the prime number theorem},
Proceedings of the Steklov Institute of Mathematics
\textbf{296} (2017), 198--210.
\href{https://doi.org/10.1134/S0081543817010163}{doi:10.1134/S0081543817010163}.

\bibitem{Revesz2025}
Sz.~Gy.~R\'ev\'esz,
\emph{The Method of Pintz for the Ingham Question About the
Connection of Distribution of $\zeta$-Zeroes and Order of the Error
in the PNT in the Beurling Context},
Michigan Mathematical Journal \textbf{75} (2025), 761--806.
\href{https://doi.org/10.1307/mmj/20226271}{doi:10.1307/mmj/20226271}.
\href{https://arxiv.org/abs/2207.00665}{arXiv:2207.00665}.

\bibitem{Broucke2026}
F.~Broucke,
\emph{On the connection between zero-free regions and the
error term in the prime number theorem},
Analysis Mathematica \textbf{52} (2026), 829--871.
\href{https://doi.org/10.1007/s10476-026-00176-y}{doi:10.1007/s10476-026-00176-y}.
\end{thebibliography}

\end{document}
